\subsection{The complete common-word quotient of the affine detector}

\begin{statement}{D23.H12}{Proved}{All affine detector markings reduce to two nine-dimensional families}
\end{statement}
Let $h=(1,f)$ be an actual normalized mixed full-GT element
whose original Hurwitz label is $(1,0)$. Its entire table of
values on $E_0^2$, where $E_0=\mathbb F_3^2\rtimes SL_2(3)$,
belongs to one of two explicitly recorded affine spaces over
$\mathbb F_3$, each of dimension nine. These are spaces of
values of one common commutator word, not independently
assigned values at different pairs. The binary tetrahedral
character of C36 is necessarily even:
\[
                          \vartheta(h)\in\{0,2\}.
                                                        \tag{D23.H12.1}
\]
Every table in these two spaces satisfies all two-cycles,
both required three-cycles at every $E_0$ pair, all pure
pentagons evaluated in $E_0$, the sixteen forced colored
values of P33, and the complete recursive raw action of
F44 on both the original and colored tuple sets. This
finite assertion does not say that every table induces an
automorphism, or that it lifts to actual GT.

\begin{proof}
\textbf{The exact common-word object.}\quad
Let $L=F(x,y)$ be the ordinary free group and let
\[
 R_E=L/\bigcap_{\psi:L\to E_0}\ker\psi.
\]
All homomorphisms to the finite group are included.
Thus this is a finite characteristic, indeed fully invariant,
quotient. C34's exhaustive rooted-Cayley inventory shows
that the $46656$ maps have exactly $146$ distinct kernels.
Choose its canonical marked pair for each kernel, and
represent $E_0$ faithfully on its nine affine points.
The two simultaneous permutations on the disjoint union
of these actions, of degree $1314$, therefore generate
exactly $R_E$.

First do the same construction for $SL_2(3)$, using the
literal linear subgroup of $E_0$. Its $576$ pairs have
$43$ distinct marked kernels from the same inventory.
Call the resulting degree-$387$ group $R_L$. The exact
orders are
\[
 |R_L|=1179648,\qquad |R_L'|=8192,
 \qquad R_L^{\mathrm{ab}}\simeq C_{12}^{\,2}.
                                                        \tag{D23.H12.2}
\]
For the last assertion, both source generators have order
twelve, so the abelianization is a quotient of
$C_{12}^{\,2}$; the computed derived index is $144$.

Projection of affine maps to their linear parts gives a
surjection $R_E\to R_L$. Its kernel $K$ consists entirely
of translations in all $146$ coordinates. Indeed a word
trivial at every linear marked pair has identity linear
part at every affine marked pair, and the converse holds.
Consequently $K$ is an elementary abelian three-group.
The exact simultaneous calculation gives
\[
 \begin{split}
 K&\simeq\mathbb F_3^{90},\qquad |R_E|=1179648\cdot3^{90},\\
 |R_E|&=10295924767167533787777935253869762453225215229952.
 \end{split}
                                                        \tag{D23.H12.3}
\]
Here is how the kernel dimension is certified. Build an
exact stabilizer chain for $R_L$, lift its strong generators
by their words in $x,y$ to the degree-$1314$ action, and
continue the chain with the old base as its initial segment.
The lifted list generates the exact joint group: the original
two generators occur in it and every other entry is their
explicit word. Strong generators fixing the old base
generate $K$. In every affine coordinate their matrices
are checked to be $I$. Row reduction of their $292$
translation coordinates over $\mathbb F_3$ has rank $90$,
and the group-order ratio is checked to equal $3^{90}$.
The saved row basis therefore gives all of $K$, without
an unverified group recognition.

Both original generators of $R_E$ still have order twelve.
Its surjection to (D23.H12.2) forces
$R_E^{\mathrm{ab}}\simeq C_{12}^{\,2}$ as well, with the
projection inducing an isomorphism on abelianizations.
In particular $K\subset R_E'$. Every lift of an element
of $R_L'$ belongs to $R_E'$.

\textbf{The four possible linear values.}\quad
Enumerate all $8192$ elements of $R_L'$. Transport their
values to all ordered pairs by the canonical marked
isomorphisms. The two-cycle at all $43$ canonical kernels
leaves $384$ elements; the positive three-cycle leaves
twelve, all of which also satisfy the negative three-cycle.
P10's normalized cubic formula forces the identity at its
linear pair, the canonical kernel $41$. Exactly four
elements remain. The known identity at kernel $37$ is
also checked and gives no further reduction.
By the literal marking of C36, these four tables have
\[
                         \vartheta=0,2,1,3,
                                                        \tag{D23.H12.4}
\]
in that order. They are called linear cases $0,1,2,3$
in the records. The complete value tables are retained,
not merely their values at $(S,R)$.

Choose a word lift $f_j\in R_E'$ of each of these four
elements, using the exact old stabilizer words. Every
element in its fiber is uniquely $k f_j$, $k\in K$.
In an affine coordinate its linear part is fixed, while
its translation is the translation of $f_j$ plus a linear
function of the $90$ coordinates of $k$. Thus these are
four affine $90$-parameter common-word families, every
member already realized by an ordinary word of $L'$.
This establishes common-word membership before imposing
any further equations.

\textbf{Literal equations on the full table.}\quad
For each of the $146$ canonical pairs, a deterministic
Cayley isomorphism transports the baseline value and every
translation-basis value to any other pair with the same
kernel. The inventory's $177$ ambient affine automorphism
orbits accelerate this transport: conjugating the two
inputs conjugates every word value. This is applied only
after exact marked-kernel identification. It produces
the complete functions on all $46656$ ordered pairs.

For fixed linear parts, a product of affine maps has
translation
\[
 v_1+A_1v_2+A_1A_2v_3+\cdots.
\]
Its dependence on the $90$ parameters is affine linear.
The same holds for inverses and conjugations of such maps.
Consequently each of the following identities is an
explicit system of linear equations over $\mathbb F_3$:
\[
 \begin{split}
 f(y,x)f(x,y)&=1,\\
 f(z,x)f(y,z)f(x,y)&=1,\\
 f(z,x)z^{-1}f(y,z)y^{-1}f(x,y)x^{-1}&=1,
                  \qquad z=(xy)^{-1}.
 \end{split}
                                                        \tag{D23.H12.5}
\]
Every linear-part equation is checked too; it is not
discarded when the translation equations are formed.
Checking the canonical kernels suffices for all pairs,
because the same marked isomorphism transports every term
of (D23.H12.5). The negative equation is necessary for
the actual element $(1,f)(-1,1)=(-1,f)$.

Append the sixteen exact values forced by P32--P33.
For each of the four linear cases, exact row reduction
leaves an affine space of dimension $15$.

We also exhaust all pure pentagons with target $E_0$.
Use H11's five-generator presentation, with images
$(a,b,c,d,e)$. For each of the $177$ ambient orbits of
the first pair, enumerate
\[
 c\in C(a),\qquad d\in C(a)\cap C(b),\qquad e\in C(b),
\]
then impose $[c,e]=1$ and
$e b^{-1}a c^{-1}b d^{-1}c e^{-1}d a^{-1}=1$.
This enumerates exactly all homomorphisms, up to only the
stated automorphism choice of the first pair. There are
$128169$ representatives. Each requires
\[
                  f(c,d)f(e,a)f(b,c)f(d,e)f(a,b)=1.
                                                        \tag{D23.H12.6}
\]
The complete equations are checked on each of the four
affine $15$-spaces. All linear parts agree and their
translation equations add no rank. Thus these whole
spaces, and not just selected points, satisfy every
$E_0$-valued pure pentagon.

\textbf{The complete raw action.}\quad
The original action is the phase-zero $V_{1,0}$ on $8064$
tuples. P33 forces colored assembly zero on $32256$ tuples.
For each four-tuple apply F44 recursively, always using
the original prefix products. The computed output must be
the literal tuple at the inverse of the stored permutation.
For fixed linear case this too is affine linear in the
fifteen remaining parameters. The exact results are
\[
\begin{array}{c|c|c|c}
\text{linear case}&\vartheta&\text{original action}&
          \text{original and colored actions}\\\hline
0&0&\text{dimension }11&\text{dimension }9\\
1&2&\text{dimension }11&\text{dimension }9\\
2&1&\text{inconsistent}&\text{inconsistent}\\
3&3&\text{inconsistent}&\text{inconsistent}.
\end{array}
                                                        \tag{D23.H12.7}
\]
All $40320$ tuples are included. A separate implementation
evaluates the sixteen affine-basis tables by literal $E_0$
multiplication, reconstructs the equations from their
differences, and obtains exactly the same affine fibers
in the original $90$ coordinates. It also checks the
independent cycle simplification
$F_k^{-1}J_k=f(P_{k-1},P_k^{-1})$, product preservation,
and the surviving outputs on every raw tuple.
This proves (D23.H12.1) and the two complete finite families.
\end{proof}

\paragraph{Exact records.}
The ordered calculation is
\path{code/sl2_3_gentle_core.py},
\path{code/affine216_universal_group.py},
\path{code/affine216_pure_pentagon_inventory.py},
\path{code/affine216_universal_constraints.py}, and
\path{code/affine216_recursive_raw_constraints.py}.
Their JSON/NPZ records contain every generator, the
translation basis, the four baseline tables, and explicit
points and nullspace bases for each surviving family.
The independent literal-action check is
\path{agents/five_level_six/affine216_prefix_recursion_replay.py}.
The whole proof and every dependency were independently read and freshly replayed by the characters worker; the five-point worker independently reconstructed the universal kernel, raw equations, and automorphism fibers. Exact check events are recorded.

\begin{statement}{D23.H13}{Proved}{Thirty-two existential normal forms for the complete affine value table}
\end{statement}
If an actual mixed full-GT element exists, then one exists
with original Hurwitz label $(1,0)$, colored assembly zero,
$\vartheta=0$, and one of the $32$ complete $E_0$ word
tables in the exact record of this statement. Of these
finite tagged residues, eight have order two, eight have
order four, and sixteen have order eight. No one of the
$32$ is asserted to have an actual lift.

\begin{proof}
Normalize as in H2 and C34. By H12 the character
$\vartheta$ is even. C36's actual Galois correction,
followed by the specified ninth power, sets it to zero
and restores the finite label $(1,0)$. P33 again forces
the colored action. Thus the entire table lies in H12's
case-zero affine nine-space, which has $3^9=19683$ elements.

An actual automorphism must induce an automorphism of
the fully invariant finite quotient $R_E$. We describe
an exact and small test for this condition. The map
\[
 (a,b)\longmapsto(a,f(a,b)^{-1}b f(a,b))
                                                        \tag{D23.H13.1}
\]
on all $E_0^2$ is bijective if and only if the induced
endomorphism of $R_E$ is an automorphism. The forward
implication follows because these homomorphisms separate
points of $R_E$: if an element of the endomorphism kernel
were nontrivial, some homomorphism to $E_0$ would detect it,
whereas surjectivity of precomposition would force every
such homomorphism to kill it. The reverse implication is
immediate. The induced endomorphism exists because the
defining intersection for $R_E$ is fully invariant.

In cases zero and one the linear value of $f$ is always
$I$ or $-I$, so (D23.H13.1) fixes each ordered pair of
linear parts $(A,B)\in SL_2(3)^2$. For a fixed pair,
the translation of a word is homogeneous linear in the
two input translations. Write its coefficient of the
second input as $K_v(A,B)$. Since the first input is fixed,
bijectivity of (D23.H13.1) on its four-dimensional
translation fiber is equivalent to
\[
           \det\bigl(I+(B-I)K_v(A,B)\bigr)\ne0.
                                                        \tag{D23.H13.2}
\]
The possible preceding scalar $I$ or $-I$ has determinant
one and does not affect this condition.

The coefficient $K_v$ is obtained by evaluating at second
input translations $e_1,e_2$ and first translation zero.
For the affine nine-parameter family each determinant is
a polynomial of degree at most two over $\mathbb F_3$.
Exhaust all $576$ linear-part pairs. Up to nonzero scalar
there are only three distinct polynomials in the stored
coordinates. Exhaustion of the $19683$ parameter vectors
leaves exactly $7776$ automorphism residues in each of the
two cases. Only case zero is needed after the actual
normalization above.

For each of these $7776$ complete case-zero tables compute
the actual GT cube and ninth power, not ordinary powers
of the word value. At a fixed pair $(a,b)$ put $b_0=b$,
and recursively
\[
 t_j=f(a,b_j),\qquad b_{j+1}=t_j^{-1}b_jt_j,
 \qquad f_n(a,b)=t_0t_1\cdots t_{n-1}.
                                                        \tag{D23.H13.3}
\]
The GT multiplication law proves that $f_n$ is the
associator value of $h^n$ whenever the original table
comes from $h$. In case zero every $t_j$ is a translation,
so the last product can also be checked by addition in
$\mathbb F_3^2$.

The exact calculation on all $146$ canonical kernels
finds that both powers keep every input in the same
allowed family, and that each has exactly $32$ distinct
outputs. For the ninth-power outputs, further applications
of (D23.H13.3) give the stated orders $2,4,8$. The output
tables, parameter vectors, and every input-to-output map
are retained. Since nine is odd, it preserves the mixed
bit; since the finite paired label is involutive, it
preserves $(1,0)$. The quartic character stays zero.
Applying this actual source operation proves the claimed
existential list of thirty-two complete affine tables.
\end{proof}

\paragraph{Certificates and limits.}
The complete determinant and tagged-power calculations are
\path{code/affine216_residue_automorphisms.py} and
\path{code/affine216_odd_power_normalization.py}.
Finite automorphisms of $R_E$ are still only necessary
residues. Compatibility with larger braid targets and an
actual profinite lift are not supplied by these calculations.

\begin{statement}{D23.H14}{Proved}{The two relative quadratic defects cannot both survive on full GT}
\end{statement}
On $\mathcal H_5=\widehat{GT}\cap\ker\lambda$, put
\[
             d_2=\Psi_3/2+a_2\pmod2.
\]
Then no element has $(\Delta,d_2)=(1,1)$. Consequently
\[
 \left.\Delta\right|_{\mathcal H_5}=0
       \quad\hbox{or}\quad
 \left.d_2\right|_{\mathcal H_5}=0.
                                                        \tag{D23.H14.1}
\]
This disjunction does not decide which character vanishes.

\begin{proof}
The supplied A.88--A.90 and E.65 make both quantities
continuous additive quadratic characters on the
multiplier-one group. On its Galois subgroup,
$\Psi_3/2$ and $a_2$ are the same quartic-root bit,
as in C36, so $d_2$ vanishes there. H2's multiplication
by an actual Galois element therefore preserves both
$\Delta$ and $d_2$ while setting $\rho\bmod3=a_2=0$.
If the initial value of $\Delta$ was one, C34 and C30
then allow an odd power with label $(1,0)$; this odd
power preserves both quadratic characters. H12 and C36
give $\Psi_3/2=\vartheta\bmod2=0$ on that element.
Since $a_2=0$, its $d_2$ is zero. Thus $(1,1)$ cannot occur.

The image of two homomorphisms to $\mathbb F_2$ is a
subgroup of $\mathbb F_2^2$. If both were nonzero, the
image would be either the diagonal subgroup, which
contains $(1,1)$, or the full group, which also contains
it. The only subgroups avoiding $(1,1)$ lie on one of
the two coordinate axes. This proves (D23.H14.1).
\end{proof}
